\documentclass[11pt]{article}
\usepackage[margin=1in]{geometry}
\usepackage{amsmath,amssymb,amsthm}
\usepackage[hidelinks]{hyperref}

\newtheorem{fact}{Fact}
\theoremstyle{definition}
\newtheorem{attack}{Attack}

\title{Disproof of TLMC Conjecture 00000000556:\\
Trivial $HH^1$ versus representation-finite rigid algebras}
\author{}
\date{}

\begin{document}
\maketitle

\begin{abstract}
Conjecture 00000000556 states that $HH^1(\Lambda)=0$ if and only if $\Lambda$
is representation-finite and rigid. We refute the ``if'' direction with the
dual numbers $\Lambda=k[x]/(x^2)$: they are representation-finite and rigid
(self-injective), yet $HH^1(\Lambda)\cong k\neq 0$.
\end{abstract}

\section{The conjecture}

\begin{quote}
$HH^1(\Lambda)=0$ if and only if $\Lambda$ is representation-finite and rigid.
\end{quote}

\noindent Here $HH^1(\Lambda)=\mathrm{Der}_k(\Lambda)/\mathrm{Inn}(\Lambda)$ is
the first Hochschild cohomology group: derivations modulo inner derivations.

\section{The counterexample}

Let $\Lambda=k[x]/(x^2)$ be the algebra of dual numbers over an arbitrary
field $k$, with basis $1,x$ and $x^2=0$.

\begin{attack}[Representation-finite]\leavevmode\\
A $\Lambda$-module is a pair $(V,T)$ where $T$ is a nilpotent endomorphism of
the vector space $V$ with $T^2=0$. Over an algebraically closed field such
pairs decompose into Jordan blocks of size at most $2$, hence the only
indecomposable $\Lambda$-modules are $k$ and $\Lambda$ itself. There are
exactly two, so $\Lambda$ is representation-finite.
\end{attack}

\begin{attack}[Rigid, i.e.\ self-injective]\leavevmode\\
The linear form $\lambda(a+bx)=b$ has associated bilinear form
$B(u,v)=\lambda(uv)$, whose matrix on the basis $\{1,x\}$ is
$\left(\begin{smallmatrix}0&1\\1&0\end{smallmatrix}\right)$ with determinant $-1$
($1$ in characteristic $2$) --- never zero. Hence $\Lambda$ is a Frobenius
algebra, therefore self-injective, i.e.\ rigid. (Independently,
$\operatorname{Ext}^1_\Lambda(\Lambda,\Lambda)=0$ trivially, since $\Lambda$
is projective as a module over itself.)
\end{attack}

\begin{attack}[$HH^1(\Lambda)\neq 0$]\leavevmode\\
Let $d\colon \Lambda\to\Lambda$ be $k$-linear with $d(1)=0$ and
$d(x)=v$. The Leibniz rule applied to $x^2=0$ forces $2xv=0$; conversely any
such $v$ yields a derivation.

\emph{Case $\operatorname{char}k\neq 2$:} $2xv=0$ forces $v\in(x)$, so
$\mathrm{Der}_k(\Lambda)\cong k$, spanned by $d=x\partial$, i.e.\
$d(1)=0$, $d(x)=x$.

\emph{Case $\operatorname{char}k=2$:} the constraint is vacuous,
$\mathrm{Der}_k(\Lambda)\cong k^2$.

Since $\Lambda$ is commutative, $uv-vu=0$ for all $u,v$, so every inner
derivation vanishes: $\mathrm{Inn}(\Lambda)=0$. The element
$d=x\partial$ satisfies $d(x)=x\neq 0$, hence is a derivation which is
\emph{not} inner. Therefore
\[
HH^1(\Lambda)\;\cong\;
\begin{cases}
k & \operatorname{char}k\neq 2,\\
k^2 & \operatorname{char}k=2,
\end{cases}
\qquad\text{in particular } HH^1(\Lambda)\neq 0.
\]
\end{attack}

\section{Conclusion}

$\Lambda=k[x]/(x^2)$ is representation-finite and rigid, but
$HH^1(\Lambda)\cong k\neq 0$. The direction
``representation-finite and rigid $\Rightarrow HH^1=0$'' of Conjecture
00000000556 fails, and with it the conjectured equivalence.

\section{Boundary of the attack}

\begin{itemize}
\item \textbf{Characteristic.} The attack works over every field: $HH^1\cong k$
  for $\operatorname{char}k\neq 2$ and $HH^1\cong k^2$ for
  $\operatorname{char}k=2$ --- nonzero either way --- and even over $\mathbb{Z}$.
\item \textbf{Meaning of ``rigid''.} The refutation adopts the
  self-injectivity reading, which the dual numbers satisfy (Frobenius).
  Under two other readings $\Lambda$ would \emph{not} be rigid and this
  particular attack would not apply:
  (i) deformation-rigidity ($HH^2=0$): the $2$-periodic $\Lambda^e$-resolution
  with differentials $x\otimes1-1\otimes x$ and $x\otimes1+1\otimes x$ gives
  $HH^2(\Lambda)\cong\Lambda/(x)\cong k\neq 0$;
  (ii) ``every indecomposable module is rigid'':
  $\operatorname{Ext}^1_\Lambda(k,k)\cong k\neq 0$ via the nonsplit extension
  $0\to(x)\to\Lambda\to k\to 0$.
\item \textbf{Failure direction.} Only the ``$\Leftarrow$'' direction is
  refuted; nothing is claimed about ``$\Rightarrow$''.
\end{itemize}

\section{Verification}

\begin{itemize}
\item \texttt{reproduce.py}: independent exact linear-algebra recomputation
  over $\mathbb{Q}$, $\mathbb{F}_2$, $\mathbb{F}_3$:
  $\dim\mathrm{Der}=1/2/1$, $\dim\mathrm{Inn}=0$, $\dim HH^1=1/2/1$,
  Frobenius determinant $\neq 0$, and $d(x)=x$ is a derivation.
\item \texttt{lean4/}: formal certificate in core Lean 4.33 (no Mathlib, no
  dependencies). \texttt{lake build \&\& lake env lean Check.lean} prints
  ``does not depend on any axioms'' for every theorem: zero axioms, zero
  \texttt{sorry}. It formalizes that $x\partial$ is a derivation, that all
  inner derivations vanish, that $HH^1\neq 0$, the classification
  $f(a,b)=(0,\,b\cdot f(x)_2)$ (so $\mathrm{Der}\cong$ base ring), and the
  nondegeneracy of the Frobenius form.
\end{itemize}

\end{document}
